
\section{Complexity bound} \label{sec:complexity}

% \sla{\bf this section is very sketchy and containts a significant gap.
% After discussions with Lukasz, 
% it seems we can bound doubly-exponentially all the steps except for
% the exhaustive search at the end of the proof of \Cref{lem:compbackfor}...
% }

In this section we provide an upper complexity bound for our decision procedure
outlined in \Cref{sec:outline}.
We will measure the size of input instances of the reachability problem 
as the size of \emph{unary}
representation of a 2-\cm $\mach$, the source $s=q(w)$ 
and the target $t=q'(w')$, namely
\[
\mysize \mach = 
\mysize Q + 2 \cdot \norm \mach \cdot \mysize {T_\text{\tiny upd}}
\ + \ \norm w \ + \ \norm w',
\]
where $T_\text{\tiny upd}$ is the set of update transitions of $\mach$.
Note that $\mysize \mach = \mysize \V$, namely we do not distinguish
between the size of $\mach$ and of the underlying 
2-\vass $\V$.
Our objective is to provide an elementary upper bound
in terms of the unary size $n=\mysize \mach$, rather than
optimizing the number of exponentials, 
and we will be content with a triply exponential bound.


\para{Bounds for \Cref{lem:targetpump}}

We estimate the values of constants $C$ and 
$D$.
The constant 
% \[
% C\ = \ C_1 (1 + e\norm\V), 
% \qquad \text{where } \quad
% C_1 \ = \ (2 C_0 e_\mathrm{max} + 1)(\ell_\mathrm{max} + 1) 2\norm\V,
% \] 
% \lukasz{
% I think it should be
\[
    C := C_1 (2 + e'_\mathrm{max}\norm\V), 
    \qquad \text{where } \quad
    C_1 = e'_\mathrm{max} 2\norm\V (\ell_\mathrm{max} + 1),
\]
% x}
is polynomial in $n$.
The norm of the effect of every cycle of length at most $e_\text{max}$ is bounded
by $e_\text{max}\norm{\V}$, and hence
all values $D_\sigma$ appearing in all the three cases
in the proof of \Cref{lem:targetpump} 
are bounded by $2(e_\text{max} \norm{\V})^2$, thus
polynomially in $n$.
Taking as $D$ the least common multiple of all $D_\sigma$ is enough for the
claim of the lemma.
Summing up:

\begin{claim} \label{claim:bounds}
  The constants $C$ and $D$ of \Cref{lem:targetpump} are bounded
  polynomially and exponentially in $n$, respectively.
\end{claim}

\para{Terminology}
For $k \in \N$ we denote $\exptriple(k) = 2^{2^{2^k}}$.
Let $X$ be a semilinear set 
given by some base-period representation.
We denote by $\norm{X}$ the maximum of norms 
and periods of this representation.
In the sequel 
we say that a semilinear set $X$ is 
\emph{(triply) exponentially semilinear} if given $\mach$,
some base-period representation of $X$ can be computed in time
(triply) exponential in $n$, 
and $\norm{X} \leq \exptriple(p(n))$ for some polynomial $p$.
%and all computed bases and periods
%are at most (triply) exponential in $n$.
Furthermore,
a mapping $S \mapsto f(S)$ transforming a semilinear set $S$
into a semilinear set $f(S)$ we call 
\emph{(triply) exponentially based} if
  given $\mach$ and a base-period representation of 
  $S\subseteq Q\times\N$, some 
  representation of $f(S)$ can be computed in time 
  (triply) exponential in $n$ and 
  $\norm{f(S)} \leq \norm{S}^d \cdot \exptriple(p(n))$
  for some polynomial $p$.
  

%  some base-period description of $f(S)$ is computable in time
%  $m$ times (triply) exponential in $n$, and the norms of 
%  all the computed bases 
%  are larger than the largest norm of a base of $S$ by at most
%  (triply) exponential multiplicative factor in $n$.


\para{Bounds for \Cref{lem:0-+}}

According to the linear-path-scheme decomposition of 
\cite[Thm.~1]{Blondin15}, whenever 
a configuration $t$ is reachable from $s$
in the underlying 2-\vass $\V$ of $\mach$,
$t$ is also reachable from $s$ by a run 
\begin{align} \label{eq:lps2}
\sigma \ \in \ 
t_0 \ c_1^{*} \ t_1 \ c_2^{*} \ \ldots \ c_\ell^{*} \ t_\ell,
\end{align}
where 
$\ell$ is polynomial in $\mysize \mach$, and
$t_0, \ldots, t_\ell$ are paths and $c_1, \ldots, c_\ell$ are cycles of
length polynomial in $\mysize \mach$ (cf.~\eqref{eq:lps}).
As shown in \cite{Blondin15},
once $\ell$, $t_0, \ldots, t_\ell$ and $c_1, \ldots, c_\ell$ are fixed, 
the reachability relation along all runs of the form \eqref{eq:lps2} 
is (the projection of) the solution set of an instance $I$ of the integer
linear programming problem of size polynomial in $\ell$, 
and hence effectively semilinear.
Moreover, numbers appearing in $I$ are bounded exponentially in
$n$, and hence by the bound of \cite{Pottier91}
(see also \cite[Prop.~4]{taming}), the norms of bases and periods
of the reachability relation along all runs of the form \eqref{eq:lps2}
are bounded exponentially in $n$.
Summing over all choices of $\ell$, $t_0, \ldots, t_\ell$ and 
$c_1, \ldots, c_\ell$, we conclude that
the reachability relation $\step^*$ of $\V$ is exponentially semilinear.
The same applies to the relation $\reachrel\ \subseteq(Q\times\N)^2$,
which is essentially the intersection of $\step^*$ with a linear set
fixing some coordinates to 0.
Finally, we transfer the same effective bounds to the three subrelations
$\stepp 0, \stepp +, \stepp -$ by noticing that each of them is
the intersection of $\reachrel$ with a semilinear set
whose bases and periods are of norm at most $C$, the constant of
\Cref{lem:targetpump}
(which is polynomial in $n$),
and intersection of linear sets in fixed dimension (here 2) 
increases the norms of bases and periods at most polynomially
(see \cite[Thm.~6]{taming}):
\begin{claim} \label{claim:cap}
Semilinear sets $r_1,r_2\subseteq (Q\times\N)^2$ satisfy
$\norm{r_1 \cap r_2} \leq (\norm{r_1} \cdot \norm{r_2})^{O(1)}$.
\end{claim}
We have thus shown:

\begin{claim} \label{claim:sl}
  The relations $\stepp 0, \stepp +, \stepp -$ are exponentially semilinear.
%   and the norms of bases and periods of their semilinear descriptions
%   are bounded exponentially in $n$.
\end{claim}

\para{Bound for \Cref{lem:0star}}

By the above effective exponential bounds on bases and periods of 
$\stepp 0$,
we conclude that the ultimately periodic sets %\eqref{eq:per1},
 \eqref{eq:per2} 
in the proof of \Cref{lem:0star} have at most exponential periods 
and bases.
In consequence, the 1-\cm $\mach'$
constructed in the proof of \Cref{lem:0star} is at most
doubly exponentially large (so is the least common multiple $N$ of
all the periods $N^c_{q,q'}$, across all choices of $c, q$ and $q'$).
This, combined with the exponential blow-up of Parikh's theorem,
yields a triple-exponential bound on the norms of bases and periods of 
the relation $\stepp 0^*$.
As both construction of $\mach'$ and Parikh's theorem are effective,
we have thus shown:

\begin{claim} \label{claim:sl2}
  The relation $\stepp{0}^*$ is triply exponetially semilinear.
%   and the norms of bases and periods of its semilinear description
%   are bounded triply-exponentially in $n$.
\end{claim}



% \para{Composition of semilinar relations}

% Before we continue the analysis we first
% have to prove that fact 
% that that composition of 
% semilinear relations increase their norms 
% polynomially.

% \begin{lemma}\label{lem:composition}
%     There is a number $d \in \N$ such that
%     for semilinear relations $r_1, r_2 \subseteq (Q\times\N)^2$
%     we have $\norm{r_1 \circ r_2} \leq O((\norm{r_1} \cdot \norm{r_2})^d)$.
% \end{lemma}
% \begin{proof}
%     We denote $r_3 = r_1 \circ r_2$.
%     \Mywlog assume $Q = \{0, \ldots, h\}$.
%     For $i = 1,2,3$ we define $r_i' \subseteq \N^2$ as
%     \[
%         r_i' := \setof{(ah + q, a'h+q') \in \N^2  }{  ((q,a), (q',a')) \in r_i }.
%     \]
%     Observe that $\norm{r_i'} = h \cdot \norm{r_i}$
%     and $r_3' = r_1' \circ r_2'$.
%     We define also 
%     \[
%         r_1'' = \setof{(a,b,i) \in \N^3}{(a,b) \in r_1', i \in \N}, \qquad    
%         r_2'' = \setof{(i,a,b) \in \N^3}{(a,b) \in r_2', i \in \N}, 
%     \]
%     and observe that $r_3'$ is obtained from $r_1'' \cap r_2''$
%     by projecting away the middle coordinate.
%     Furthermore,
%     $\norm{r_1''} = \norm{r_1'}$, $\norm{r_2''} = \norm{r_2'}$.
%     We conclude from \cite[Thm.~6]{taming} that there is a polynomial 
%     $p$ independent of sets $r_1,r_2$, such that 
%     $\norm{r_3'} \leq p(\norm{r_1''}, \norm{r_2''})$.
%     Hence, the lemma follows.
% \end{proof}


\para{Bound for \Cref{lem:compbackfor}}

Composition of two semilinear relations in fixed dimension, 
very much like intersection,
increases the norms of bases and periods at most polynomially:

\begin{claim} \label{claim:circ}
Semilinear relations $r_1,r_2\subseteq (Q\times\N)^2$ satisfy
$\norm{r_1 \circ r_2} \leq (\norm{r_1} \cdot \norm{r_2})^{O(1)}$.
\end{claim}
Indeed, ignoring $Q$ fro simplicity, \Cref{claim:circ} 
follows immediately from \cite[Thm.~6]{taming}, similarly to 
\Cref{claim:cap}, as
the composition $r_1 \circ r_2$ is the projection to the first and last coordinate
of 
\[
(r_1 \times \N) \ \cap \ (\N\times r_2).
\]
Using \Cref{claim:sl2} and \Cref{claim:circ} we conclude
that the relation
  $(\stepp - \circ \, \stepp{0}^*)$ is triply exponentially semilinear,
which allows us further to deduce the bounds on bases
of the (necessarily $D$-saturated) inverse image of the relation,
as well as on the bases of the semilinear image thereof:

\begin{claim}\label{claim:sl3}
The mapping $S \mapsto (\stepp - \circ \, \stepp{0}^*) \,\circ\, S$
  is triply exponentially based.
\end{claim}

\begin{claim}\label{claim:sl3a}
    The mapping
  set  $S\mapsto S \,\circ\,(\stepp - \circ \, \stepp{0}^*)$
  is triply exponentially based.
\end{claim}


Concerning the first part of \Cref{lem:compbackfor},
we bound the number of iterations of the transformation
\begin{align}\label{eq:tra}
S \quad \mapsto \quad S \ \cup \ 
(\stepp - \, \circ \, \stepp{0}^*)  \circ  S
\end{align}
needed to compute the fixpoint $\Phi_3 \circ \{t\}$, starting from
$X_0 = ({\stepp -} \,\circ\ {\stepp{0}^\ast})  \circ \{t\}$.
Due to \Cref{claim:sl3}, the bases of the $D$-saturated 
set $X_0$ are bounded triply exponentially in $n$.
Then every base of every next $D$-saturated set
\[
    X_{i+1} \ := \ X_{i} \ \cup \ 
    ({\stepp -} \,\circ\, \stepp{0}^\ast) \, \circ\,  X_{i},
\]
as $X_i \subseteq X_{i+1}$,
is either smaller than some base of $X_i$, or is 
\emph{newly created}, namely
 $\minf{X_{i+1}}(q,r) = m < \infty$ while $\minf{X_i}(q,r)=\infty$
 for some $q\in Q$ and $r\in\setfromto 0 {D-1}$.
 A newly created base may appear at most $\mysize Q \cdot D$ 
 times (exponentially many times) during the whole iteration.
 Due to \Cref{claim:sl3}
 there is $d \in \N$ and polynomial $p$ such that if new base appears in $X_{i+1}$ then
 \[
    \norm{X_{i+1}} \leq \norm{X_i}^{d} \cdot \exptriple(p(n)),
 \]
 \lukasz{Probably it would be possible to prove that $X_{i+1}$ is only multiplied 
 by some triply exponential factor, without the exponent $d$.
 However, at the moment I don't see how, and the complexity is the same anyway.}
 otherwise $\norm{X_{i+1}} \leq \max\{\norm{X_i}, D\}$.
 %larger than a base of $X_i$ by
 %at most triply-exponential multiplicative factor, due to 
 %\Cref{claim:sl3}.
In consequence for every $X_i$ we 
have $\norm{X_i} \leq \norm{X_0}^{d \cdot \mysize Q \cdot D} \exptriple(p'(n))$
for some polynomial $p'$,
which yields a triply exponential bound on $\norm{X_i}$.
Therefore, the total number of iterations of \eqref{eq:tra} 
before 
a fixpoint is reached is at most triply exponential as well.
% Multiplying a triple exponential value exponentially many times
% yields still a triple exponential quantity,  
%therefore all newly created bases are bounded triply exponentially.
%In consequence, the total number of iterations of \eqref{eq:tra} 
%before 
%a fixpoint is reached is at most triply exponential as well.
Our reasoning was not specific to the starting set $\{t\}$,
and we have actually shown:

\begin{claim}\label{claim:sl4}
    The mapping $S\mapsto \Phi_3 \circ S$
  is triply exponentially based and computable in
  at most triply exponentially many iterations of \eqref{eq:tra}.
\end{claim}

% \smallskip

Concerning the second part of \Cref{lem:compbackfor},
we need to bound the number of iterations 
of the transformation
\begin{align}\label{eq:tra2}
S \quad \mapsto \quad S \ \cup \ f(S),
\qquad \text{ where } \quad
f(S) \ = \ S \ 
\circ \ \bigl(\stepp{0}^* \ \circ \ \Phi_3 \ \circ \ \stepp +\bigr),
\end{align}
needed to compute the
fixpoint $\{s\} \circ \Phi_1$, starting from
$X_0 =  f\bigl(\{s\}\bigr)$. 
We should also provide an upper bound on computation of $X_0$,
and computation of every iteration
$
X_{i+1}  :=  X_i  \cup f\bigl(X_i\bigr).
$

\medskip

We start with the computation of $X_0$.
By \Cref{claim:sl4} applied to $S= (\stepp + \circ \ X_{q,r})$,
the $D$-saturated set $(\Phi_3 \circ \stepp +)\circ X_{q,r}$
is computable in at most triply exponentially many iterations
of the transformation \eqref{eq:tra}, and its intersection $Y$
with the set $(S \, \circ \stepp{0}^*)$, 
\[
Y \ = \ (S \, \circ \stepp{0}^*) \ \cap \ 
\bigl((\Phi_3 \, \circ \stepp +)\circ X_{q,r}\bigr)
\]
(relying on \Cref{claim:sl2} and on polynomial blowup of intersection
of semilinear sets in fixed dimension \cite[Thm.~6]{taming})
has bases and periods of at most triply exponential norm. 
Now, a key observation is that since  the set
$(\Phi_3 \circ \stepp +)\circ X_{q,r}$ is computed in at most
triply exponentially many (say $m$) iterations of \eqref{eq:tra},
the set 
\[ Y \ \circ \ (\stepp - \, \circ \, \stepp{0}^*)^m
\]
has necessarily nonempty intersection with $X_{q,r}$, and moreover
the intersection contains, due to \Cref{claim:sl3a}, a number
which is at most triply exponential in $n$.
This provides a triply exponential bound on the number of iterations 
of the exhaustive search 
in the proof of \Cref{lem:compbackfor}, and thus the bound
on the computing time for $X_0$.
In the same way one obtains the bound on the computing time for every
next iterate $X_{i+1}$.

Finally, the triple exponential bound on the number of iterations 
\eqref{eq:tra2} before a fixpoint is reached in show similarly
as in the first part.
This completes the complexity analysis of the second part of 
\Cref{lem:compbackfor}.

\medskip

% Applying the pushdown automata saturation procedure of \cite{BEM97} 
% to $\mach'$, we conclude that
% the bases and period of $\stepp 0 \circ S$ and $S \circ \stepp 0$
% can be larger than 
% the corresponding quantities of $S$ by at a value at most 
% doubly exponential in $n$.
% Also, composition with $\stepp -$ or $\stepp +$ cannot increase
% bases or periods by larger quantity, due to the exponential bounds
% on bases and periods of semilinear descriptions given by combination
% of \cite{Blondin15} and \cite{Pottier91}.
% We apply these observations to bound the bases of $D$-\saturated sets 
% $X_i$ arising in the iterative computation in the proof of 
% \Cref{lem:compbackfor}.
% In every iteration, either a base is lowered, or a new base appears,
% namely $\minf{X_{i+1}}(q,r) = m < \infty$ while $\minf{X_i}(q,r)=\infty$.
% Whenever this happens, due to the above observations
% the new value $m$ is larger
% than the maximal base in $X_i$ by at most doubly exponential value.
% As this happens at most $\mysize Q D$ times, 
% the largest base ever appearing in any of $X_i$ is at most
% doubly exponential.

\medskip

Summing up, the running time of our decision procedure is at
most triply exponential in $n$, the size of unary representation,
and hence at most quadruply exponential in the size of binary 
representation,
while the only known lower bound is \pspace-hardness
(in binary representation).
Narrowing the apparent complexity gap is left as future research.
